\scanpage{148}
\mspage{73}
sans hyp. superflues, \unsure{je} vais plutôt (comme dans Mar p. \unsure{779|769|709} ff) définir une flèche en sens inverse
\[
(243)\qquad\Pi_1^{\mathrm{arcs}}P_{D'^{*}}\longrightarrow\Pi_1^{\mathrm{rev}}\mathcal V^{*}_{D',Y},
\]
où le premier $\Pi_1$ est défini par pts et arcs, et le deuxième (on n'a pas le choix !) est défini en termes de \unsure{rev. + fonc. can.}. Il faut donc généraliser en contexte présent les considérations du loc. cit. Je ne vais pas entrer dans le détail ici --- d'autant moins que je ne vois pas clairement à présent dans quelle mesure le th. \inserted{envisagé} est indispensable --- s'il n'est pas tout aussi commode \deleted{\unsure{et utile}} \inserted{pour la suite} de travailler avec les $\Pi_1\mathcal V^{*}_{D',Y}$ au lieu des $\Pi_1P_{D'^{*}}$ --- en modifiant dans ce sens les définitions des $\widetilde\Pi_{d_0,d_1,\ldots,d_r}$ ; Bien entendu, on aura besoin par la suite que $\widetilde\Pi_{d_0,\ldots,d_r}\mathbin{-}\Pi_{d_0,\ldots,d_r}$ ait une \uline{structure d'extension}, de noyau un quotient de $T_{d_0,\ldots,d_r}$ etc.
\margin{147}
